\clearpage
\VerdierPhysicalPageAnchor{0266}%
\begingroup
\small
\setlength{\parindent}{0pt}%
\setlength{\parskip}{0pt}%
\setlength{\emergencystretch}{2em}%
\newcommand{\VerdierBibRow}[3]{%
  \VerdierSemanticAnchor{#1}%
  \par\noindent\hangindent=3em\hangafter=1
  \makebox[3em][l]{[#2]}#3\par\smallskip}

\vspace*{4.5\baselineskip}
\VerdierSemanticAnchor{bib-title}%
\begin{center}
  {\Large\bfseries 参考文献}
\end{center}
\vspace{2.5\baselineskip}

\VerdierBibRow{bib-001}{1}{H. Cartan と S. Eilenberg,「ホモロジー代数」,
  Princeton University Press (1956).}

\VerdierBibRow{bib-002}{2}{D. Puppe, \emph{安定ホモトピー論の形式構造について},
  Colloquium on Algebraic Topology, Aarhus Universitet, pp. 65-71 (1962).}

\VerdierBibRow{bib-003}{3}{N. Yoneda, \emph{\textup{Ext} と完全列について},
  J. Fac. Sci. Univ. Tokyo. Sec. I, Vol. VIII, pp. 507-576 (1960).}

% Source-fidelity note: the printed authority reads "Harsthorne".
\VerdierBibRow{bib-004}{4}{R. Hartshorne,「留数と双対性」, Lecture Notes in
  Mathematics n° 20, Springer-Verlag (1966).}

\VerdierBibRow{bib-005}{5}{J.-L. Verdier,
  \emph{スキームのエタール・コホモロジーにおける双対定理},
  Lecture Notes of the Woods Hole Conference,
  American Mathematical Society (1964).}

\VerdierBibRow{bib-006}{6}{M. Artin、A. Grothendieck、J.-L. Verdier,
  「トポスの理論とスキームのエタール・コホモロジー」 (S.G.A. 4),
  Séminaire de Géométrie Algébrique du Bois-Marie (1963-1964),
  Lecture Notes in Mathematics n° 269, Springer-Verlag (1972).}

\VerdierBibRow{bib-007}{7}{J.-P. Serre,「Galois コホモロジー」,
  Lecture Notes in Mathematics n° 5, Springer-Verlag (1964).}

\VerdierBibRow{bib-008}{8}{M. Artin と J.-L. Verdier,
  「数体のエタール・コホモロジーに関するセミナー」,
  Lecture Notes of the Woods Hole Conference,
  American Mathematical Society (1964).}

\VerdierBibRow{bib-009}{9}{J.-L. Verdier,
  \emph{局所コンパクト空間のコホモロジーにおける双対性},
  Séminaire Bourbaki, exposé n° 300 (1965-1966).}

\VerdierBibRow{bib-010}{10}{J.-L. Verdier と M. Zisman,
  \emph{Poincaré 双対性}, Séminaire Heidelberg-Strasbourg (1966-1967),
  Publ. I.R.M.A. Strasbourg n° 3 (1969).}
\endgroup
